\section{Examples and remarks}\label{Arden24-sec-4}

\begin{remark}[{\cite[Remark 4.1, p.~2]{Arden24}}]\label{Arden24-rem-4.1}
\label{Arden24-rem:orientation}
Nothing above uses the orientation of the arrows except as bookkeeping, which is the content of the last sentence of the proof of Theorem~\ref{Arden24-thm-3.1}. Readers who prefer graphs to quivers may replace every quiver by its underlying graph with an arbitrary orientation, and no statement changes.
\end{remark}

\begin{remark}[{\cite[Remark 4.2, p.~2]{Arden24}}]\label{Arden24-rem-4.2}
\label{Arden24-rem:infinite}
For an infinite quiver of finite valence the divergence map is still defined on finitely supported weightings and its kernel is still saturated, but the rank formula of Theorem~\ref{Arden24-thm-3.1} has no meaning, since all three of its terms are infinite. The correct statement in that setting is a description of the cycle space as a direct limit, which we do not pursue.
\end{remark}
